% 2026 narrative draft; source mapping and open items are in planning/2026-draft-status.md.
% Shared presentation remains in shared/preamble.tex.
\firstname{David}
\familyname{Molik}
\title{Directed Service -- 2026}
\address{University Library}{Kansas State University}
\portfoliotitle

\section{OPE Department Head and Lead Research Officer}
My assigned leadership responsibilities connect open publishing, research-data services and institutional research capacity. The Lead Research Officer MOU describes proposal and partnership development, collaborative funding, project scoping, writing and budgeting. The May 1 administrative correspondence confirms execution of the agreement, and August 5 Strategic Leadership Council minutes corroborate the role and my attendance. This complements my continuing OPE Department Head appointment.\evidenceref{L1}

\subsection*{Supervision, delegation and mentoring}
I developed three student roles spanning data/digital scholarship, publishing/repositories and OER/accessibility. The July--August correspondence records my rationale for early recruitment, meaningful technical work, a primary librarian mentor and increasing responsibility. I revised the descriptions, obtained approval, and sent them to HR. HR confirmed all three Handshake postings and supervisor ownership on August 10. September Teams records document the addition and welcome of three students. I proposed a paid initial open-science training pathway and authorized supervisors to use it; completion and learning outcomes are not yet documented.\evidenceref{L2}

I supported faculty research development through iterative mentoring rather than a single consultation. In July, I reviewed Gwen Sibley's William T. Grant letter of inquiry, focusing on research questions, the distinction between data and research evidence, methods, sponsor alignment and biosketch presentation. Her July 22--23 replies describe incorporating my recommendations, and I circulated Draft 8 on July 24. This provides direct evidence of advice used in revision while preserving her authorship and project ownership.\evidenceref{L3}

I also delegated OPE representation for an Academic Services session and secured access to Carpentries instructor-training opportunities for colleagues. The correspondence establishes coordination and access, not subsequent presentation delivery or certification.\evidenceref{L2}\evidenceref{D1}

\subsection*{Developing research partnerships and services}
Following an August 11 discussion with Jeff Sheldon, I delivered a Sunderland operating-model workbook to Jeff, Tara Coleman and James Wiser on August 12. The accompanying analysis connected ATE/IUSE opportunities to a sustainable research-support program and distinguished routine services from advanced technical and research responsibilities. Jeff's response recognized the write-up. By September 11, we had deferred the ATE application to the following year to strengthen the MATC relationship and other partnerships. The delivered model is the accomplishment; implementation and a funded program remain prospective.\evidenceref{L4}

On September 1, I met with Carol Sevin about grant development and shared USDA/NSF capacity-building resources in the meeting chat. Separately, Carol confirmed incorporating my ORCID suggestion into research guidance on September 2. With melia fritch, Gwen Sibley and Jeff Sheldon, I began developing an IUSE concept examining maker-based learning, STEM self-efficacy and transfer across disciplines. I identified the opportunity, formulated a comparative question, recruited collaborators and shared an example before the September 16 planning meeting. This draft reports preparation, not outcomes of that future meeting.\evidenceref{L4}

For Cassandra Olds's SAS concept, I developed an OPE data-infrastructure scope with an initial technical leadership role for me and a transition to Gwen Sibley's operational RDM leadership. The proposed five-year OPE component totals \$448,798; this is neither the full project budget nor an award. I also developed proposed open-research/publishing services for the OAK-AI Hub discussion. These activities demonstrate service scoping and partnership development, with sponsor and implementation outcomes still separate.\evidenceref{G1}

\subsection*{Research-publication reporting}
The June 1 OVPR request asked Libraries to help assess publication indicators and support NextGen reporting. My response addressed citation windows, affiliations, source quality, open-access categories and research outputs beyond articles. I coordinated Libraries participation and developed complementary infrastructure and reporting approaches. HAIL separates normalized, provenance-aware collection from presentation; HAIL-local provides a cached reporting route; Prairie Observatory explored lightweight repository/dashboard views. These are related efforts serving one institutional need, not three independent impact claims.\evidenceref{K1}

August--September correspondence and reports establish preliminary output generation and review. Hans Coetzee praised development progress on August 27, and Lowell Scott provided positive visual feedback in September. Report versions differ, however, and an August 31 fetching log was followed by a failure notice for that job. I therefore report prototype development, design decisions and preliminary delivery, while reserving claims about reliable full refreshes, final institutional metrics and decision use for validated evidence. Namratha Kommineni's Power BI prototype is a collaborator's contribution, not my sole work.\evidenceref{K1}

\subsection*{Publishing, data stewardship and outreach}
I coordinated service-development discussions following the reported K-REx DSpace 9.3 upgrade, addressing ORCID, OpenAIRE, Crossref, preservation and research-reporting interfaces. This is coordination and planning evidence, not a claim that I personally performed the upgrade or completed every integration. My June Research Weekly article, ``What is the Serials Crisis?'', translated publishing conditions for the campus community; leadership correspondence confirms publication and favorable qualitative response. I provided participants with follow-up research-data and publishing resources after the May 29 workshop.\evidenceref{D1}

The achievements record also documents Textbook 2.0 receipts of \$1,920 against \$385 in 2025, and a consultation involving approximately 250 TB of research data. These figures describe the reported fundraising comparison and consultation scale; they do not establish student savings or data deposited. An identified Green OA route for two chapters similarly remains a possible cost-avoidance opportunity, not realized savings.\evidenceref{D2}
